\section{What direct substitutions into existing bounds establish}
\label{sec:prioranalysis}
This comparison concerns the stated bounds below, rather than every
consequence of the cited papers. It does not establish bibliographic priority.

Pedregosa et al.~\cite{pedregosa}, Theorem 3, give a convex primal--dual
bound in terms of averaged accepted smoothness estimates. Thus adaptive
curvature and last-iterate dual guarantees are already present there.
One cannot identify that dual value with the maximum of scalar tangent
lower bounds used here: their dual variable averages gradients. This
distinction alone does not prove a new rate; a transfer argument along
the algorithm's trajectory could still exist.

Here is an elementary illustration of the certificate distinction. For
$f(x)=x^2/2$ on $[-1,1]$, the scalar tangent lower bound is
$a(x)=-x^2/2-|x|$, whereas the Fenchel dual value is
$\psi(u)=-u^2/2-|u|$. For the history $(-1,1)$,
\[
 \max_i a(x_i)=-3/2 < \psi((x_0+x_1)/2)=0.
\]
For the history $(-1,-1/2)$,
\[
 \max_i a(x_i)=-5/8 > \psi(-3/4)=-33/32.
\]
Thus the two lower bounds have no universal ordering for arbitrary
tangent histories. These histories are not asserted to be iterates of
the cited algorithm; neither example contradicts its theorem.

A global Euclidean-smoothness substitution also loses information.
Let $\Delta_Q=\max_{x\in Q}(f(x)-f_Q^*)$. Assumption~\eqref{eq:curvature}
gives a valid bound $L_Q=H_0+H_1\Delta_Q$ on $Q$. Applying a conventional
$O(L_QD^2/\varepsilon)$ estimate retains $\Delta_Q/\varepsilon$, rather
than the logarithmic $H_1$ term. Replacing $\Delta_Q$ by $U_0$ is not
automatically valid: $U_0$ only bounds the initial suboptimality.
Restricting $Q$ to an initial sublevel set changes its LMO. Our local
acceptance argument avoids that restriction, followed by a phase-wise
count of estimates that remain too large. A finer use of an existing
average-curvature bound remains an open priority question.

The relative-smooth framework of Takahashi, Pokutta, and
Takeda~\cite{bregmanfw} uses a kernel, a segment scaling condition, and
additional assumptions for its accelerated conclusions. Its Euclidean
specialization recovers ordinary smoothness bounds, so choosing a
Euclidean kernel and $L=L_Q$ alone does not yield the separated
$H_0,H_1$ dependence above. This does not exclude a more suitable
kernel or a refined specialization of that analysis.

For adaptive sliding~\cite{adcgs}, Corollary 3.3 concerns outer
iterations with a running maximum of local curvature estimates;
Theorem 2.1 separately bounds inner LMO effort. Its restarted
Theorem 4.1 additionally assumes local strong convexity and uses local
smoothness and strong-convexity parameters in the restart schedule.
Our curvature assumption does not imply strong convexity: for example,
$f(x_1,x_2)=\cosh(x_1)-1$ on $[-1,1]^2$ has $H_0=H_1=1$ and a line
of minimizers. Consequently the restarted strong-convexity theorem
cannot be invoked from our assumptions alone. A complete comparison
must count all inner LMO calls and preserve the stopping certificate;
we do not establish superiority over sliding methods.

Wirth, Pe\~na, and Pokutta~\cite{wirth}, Definition 2.3(b) and
Theorem 3.6, use a uniform extended-curvature constant and a vertex
error bound for exact line search or a short step. Their linear
and two-regime rates are prior results. Directly substituting a
global $L_QD^2$ retains global curvature; a varying-curvature extension
would require further analysis.

Kerdreux, d'Aspremont, and Pokutta~\cite{restartfw}, Proposition 4.2
and Theorem 5.1, already combine burn-in, computable strong-Wolfe
stopping, unknown error-bound parameters, and cumulative LMO counts.
The stated bound uses global away curvature. Our phases analyze
ordinary backtracking without algorithmic restarts; this difference
does not itself establish a novel bound.

Zhao~\cite{zhao}, Theorems 4.1--4.2, establishes linear rates for both
objective error and FW gap with logarithmically homogeneous barriers.
Algorithm 1
uses exact search or a local-Hessian analytic step, rather than our
first-order quadratic backtracking. The function assumptions differ.
The passage from primal to gap rates is therefore also prior in
spirit, and is not claimed as novel here.
